% !TeX program = xelatex
\documentclass[12pt,a4paper]{article}
\usepackage{fontspec}
\setmainfont{Liberation Serif}   % или Times New Roman, Arial
\usepackage[english,russian]{babel}   % русский и английский

\usepackage{amsmath,amssymb,amsthm,mathtools}
\usepackage{geometry}
\usepackage{booktabs}
\usepackage{hyperref}
\usepackage{url}
\usepackage{xcolor}
\usepackage{enumitem}
\usepackage{tikz}
\usetikzlibrary{decorations.pathreplacing,arrows.meta,positioning,shapes.geometric}
\usepackage{float}

% Макросы YUCT
\newcommand{\Keff}{K_{\mathrm{eff}}}
\newcommand{\kappac}{\kappa_c}
\newcommand{\eps}{\varepsilon}
\newcommand{\df}{d_{\!f}}

\hypersetup{colorlinks=true,linkcolor=blue,citecolor=blue,urlcolor=blue}

\begin{document}
	
	\begin{titlepage}
		\centering
		\vspace*{1cm}
		{\Huge\bfseries Приложение S2: Фрактальная размерность координации в двумерных гетероструктурах Шоттки и \\ предсказание показателя масштабирования YUCT $\beta = 2/3$ ($\approx 0.67$)}\\[1cm]
		{\large Алексей В. Якушев}\\[0.3cm]
		{\normalsize Исследования Якушева, \url{https://yuct.org/}}\\[1.5cm]
		{\normalsize DOI: \texttt{10.5281/zenodo.18444598}}\\[0.5cm]
		{\normalsize Май 2026\\Версия 1.0}\\[2cm]
		
		\begin{minipage}{0.9\textwidth}
			\begin{abstract}
				\noindent
				Ang, Yang и Ang (Phys. Rev. Lett. 121, 056802, 2018) обнаружили
				универсальные законы масштабирования тока от температуры в двумерных
				гетероструктурах Шоттки: $\log(\mathcal{I}/T^{3/2}) \propto -1/T$ для
				поперечных контактов и $\log(\mathcal{I}/T^{1}) \propto -1/T$ для
				вертикальных контактов с несохраняющимся импульсом. В этом приложении
				мы показываем, что оба показателя определяются эффективной фрактальной
				размерностью $d_f$ координационной сети, управляющей переносом.
				Универсальный закон масштабирования ошибок YUCT $\varepsilon = \kappa_c\alpha
				K_{\mathrm{eff}}^{-\beta}$ с $\beta = 2/3 \approx 0.67$ связан с $d_f$
				через $\beta = 2/d_f$. Из экспериментального масштабирования тока мы
				извлекаем $d_f = 3$ для поперечных контактов и $d_f = 2$ для
				вертикальных контактов, что даёт $\beta = 2/3$ и $\beta = 1$
				соответственно. Анализ предсказывает, что относительные флуктуации тока
				в поперечных диодах Шоттки должны следовать закону
				$\varepsilon_{\mathrm{fluc}} \propto K_{\mathrm{eff}}^{-2/3}$ —
				проверяемое предсказание, которое можно проверить с помощью стандартной
				шумовой спектроскопии. Эта работа включает двумерные гетероструктуры
				Шоттки в семейство систем, которые демонстрируют или, как ожидается,
				будут демонстрировать универсальный показатель YUCT $\beta = 2/3$
				($0.67$), расширяя область применимости теории на твердотельную
				наноэлектронику.
			\end{abstract}
			\vspace{1cm}
			\noindent\small\textbf{Ключевые слова:} YUCT, эффективность координации, универсальное масштабирование,
			$\beta = 2/3$, 0.67, гетероструктуры Шоттки, 2D-материалы,
			фрактальная размерность, универсальность КПЗ.
		\end{minipage}
	\end{titlepage}
	
	\tableofcontents
	\newpage
	
	\section{Введение}
	Объединяющая теория координации Якушева (YUCT) постулирует, что любой
	процесс координации подчиняется универсальному закону масштабирования ошибок
	\begin{equation}
		\varepsilon = \kappa_c \, \alpha \, \Keff^{-\beta},
		\qquad \beta = 2/3 \approx 0.67,\;\; \kappa_c = 1/3,
		\label{eq:yuct-law}
	\end{equation}
	где $\Keff$ — эффективность координации, а $\alpha$ — системно-зависимый
	предэкспоненциальный множитель. Этот закон был проверен на системах
	от химических связей до космологических флуктуаций
	(Приложения L, M, E, F, G, V).
	
	Ang, Yang и Ang~\cite{Ang2018} вывели универсальные законы масштабирования
	тока от температуры для гетероструктур Шоттки на основе 2D-материалов:
	\begin{align}
		\text{поперечные контакты:}&\quad \log\Bigl(\frac{\mathcal{I}}{T^{3/2}}\Bigr) \propto -\frac{1}{T},
		\label{eq:ang-lat}\\
		\text{вертикальные контакты ($k_\parallel$ не сохраняется):}&\quad \log\Bigl(\frac{\mathcal{I}}{T^{1}}\Bigr) \propto -\frac{1}{T}.
		\label{eq:ang-vert}
	\end{align}
	Показатели $\beta_{\mathrm{IV}} = 3/2$ и $1$ не зависят от конкретного
	2D-материала и отражают геометрию пути тока.
	
	В этом приложении мы демонстрируем, что показатель $\beta_{\mathrm{IV}}$
	не является независимым числом, а определяется эффективной фрактальной
	размерностью $d_f$ координационной сети, управляющей переносом.
	Из экспериментального $\beta_{\mathrm{IV}}$ мы извлекаем $d_f$, а из $d_f$
	предсказываем показатель ошибки координации $\beta_{\mathrm{err}} = 2/d_f$.
	Для поперечных контактов это даёт $\beta_{\mathrm{err}} = 2/3$ — в точности
	универсальное значение YUCT. Масштабирование флуктуаций, вытекающее из
	этого предсказания, представляет собой проверяемый тест теории.
	
	\section{YPSDC-словарь для барьера Шоттки}
	На языке YPSDC (Протокол синхронной распределённой координации Якушева)
	поперечная гетероструктура Шоттки является координационной системой со
	следующими компонентами:
	\begin{itemize}
		\item \textbf{Словарь $\mathcal{D}$:} множество одноэлектронных
		состояний в 2D-плоскости с энергией $\epsilon > \Phi_B$ и
		волновым вектором $\mathbf{k}$, удовлетворяющим внутриплоскостному
		закону дисперсии (\ref{eq:dispersion}).
		\item \textbf{Индекс $\kappa$:} тепловая энергия $k_B T$, которая
		переводит электрон в словарь.
		\item \textbf{Резонанс $R$:} сохранение тангенциального импульса
		$k_y$ при прохождении через барьер.
	\end{itemize}
	Координированным действием является успешная термоэлектронная эмиссия
	электрона, который одновременно удовлетворяет условию по энергии
	$\epsilon > \Phi_B$ и условию по импульсу $k_x > 0$ (направлен к барьеру).
	Эффективность этой координации определяет обратный ток насыщения.
	
	\section{Экспериментальное масштабирование тока и эффективности координации}
	
	Для общей изотропной 2D-дисперсии
	\begin{equation}
		\epsilon(\mathbf{k}) = \sum_n c_n |\mathbf{k}|^n,
		\label{eq:dispersion}
	\end{equation}
	Ang и др. показали, что число активных каналов (способных вносить вклад в ток)
	универсально масштабируется как
	\begin{equation}
		N_{\mathrm{ch}} \propto T^{1/2}.
		\label{eq:Nch}
	\end{equation}
	Полное число термически доступных состояний выше барьера пропорционально
	$T$ (из-за хвоста распределения Больцмана). Следовательно, доля успешно
	скоординированных состояний
	\begin{equation}
		\Keff(T) = \frac{N_{\mathrm{ch}}}{N_{\mathrm{tot}}} \propto T^{-1/2},
		\label{eq:keff-T}
	\end{equation}
	служит эффективностью координации для данного процесса. Измеряемый ток
	принимает вид
	\begin{equation}
		\mathcal{I} = \mathcal{I}_0 \, \Keff(T)^{-3} \, e^{-\Phi_B/k_B T},
		\label{eq:I-via-keff}
	\end{equation}
	что воспроизводит экспериментальный закон (\ref{eq:ang-lat}).
	
	Операциональное определение $\Keff$, пригодное для шумовых измерений,
	имеет вид
	\begin{equation}
		\Keff = \frac{I}{I_0},
		\label{eq:keff-operational}
	\end{equation}
	где $I$ — измеренный ток, а $I_0$ — эталонный ток, который тек бы при
	отсутствии координационного ограничения. $I_0$ может быть получен
	экстраполяцией высокотемпературного поведения вертикального контакта
	(где асимптотически $\Keff(T) \propto T^{-0}$) к температуре измерения,
	либо с помощью независимой калибровки по 3D-диоду Шоттки с той же высотой
	барьера. Это операциональное определение позволяет определять $\Keff$
	непосредственно из вольт-амперных характеристик, не полагаясь на
	теоретическое масштабирование (\ref{eq:keff-T}).
	
	\section{Фрактальная размерность координационной сети}
	
	Центральный результат YUCT (Приложение L) состоит в том, что показатель
	масштабирования ошибок определяется фрактальной размерностью $d_f$
	координационной сети:
	\begin{equation}
		\beta_{\mathrm{err}} = \frac{2}{d_f}.
		\label{eq:beta-df}
	\end{equation}
	Плотность активных координационных каналов в тепловом окне несёт
	геометрическую сигнатуру сети. В изотропном $d_f$-мерном пространстве
	плотность состояний, доступных при температуре $T$, масштабируется как
	$T^{d_f/2 - 1}$. Для поперечного барьера Шоттки дополнительное кинематическое
	ограничение, связанное с сохранением импульса (множитель $\cos\theta$),
	уменьшает эффективный показатель на единицу, так что кинетический показатель
	становится
	\begin{equation}
		\beta_{\mathrm{IV}} = \frac{d_f}{2}.
		\label{eq:bIV-df}
	\end{equation}
	\textbf{Уравнение (\ref{eq:bIV-df}) является гипотезой, мотивированной
		геометрическим сходством между путём тока и координационной сетью с
		фрактальной размерностью $d_f$. Строгий вывод из микроскопического
		действия YUCT пока не проведён; поэтому соотношение следует рассматривать
		как предположение, приводящее к проверяемым предсказаниям (раздел \ref{sec:predictions})
		и приглашающее к дальнейшим теоретическим исследованиям.}
	
	Уравнения (\ref{eq:beta-df}) и (\ref{eq:bIV-df}) объединяются в фундаментальное
	соотношение
	\begin{equation}
		\boxed{\beta_{\mathrm{IV}} \cdot \beta_{\mathrm{err}} = 1}.
		\label{eq:relation}
	\end{equation}
	
	Из экспериментального значения $\beta_{\mathrm{IV}} = 3/2$ получаем
	\begin{equation}
		d_f = 2\beta_{\mathrm{IV}} = 3, \qquad
		\beta_{\mathrm{err}} = \frac{2}{d_f} = \frac{2}{3} \approx 0.67.
	\end{equation}
	Таким образом, измеренное масштабирование тока определяет координационную
	сеть с эффективной фрактальной размерностью $d_f = 3$, что, в свою очередь,
	означает, что любая мера координационной ошибки (например, относительные
	флуктуации тока) должна подчиняться степенному закону YUCT с показателем
	$2/3$. Это \textbf{предсказание}, а не подгонка под уже известные данные,
	поскольку измерения флуктуаций на этих системах ещё не проводились.
	
	Для вертикальных контактов с несохраняющимся $k_\parallel$ угловое
	ограничение исчезает, $\beta_{\mathrm{IV}} = 1$, и, следовательно,
	$d_f = 2$, $\beta_{\mathrm{err}} = 1$. YUCT предсказывает, что флуктуации
	тока в таких контактах будут масштабироваться как $\Keff^{-1}$, в отличие
	от закона $2/3$.
	
	\section{Предсказания для флуктуаций тока}
	\label{sec:predictions}
	Применение закона ошибок YUCT (\ref{eq:yuct-law}) к эффективности
	координации (\ref{eq:keff-T}) даёт количественное предсказание для
	мощности низкочастотного шума тока $S_I$ в поперечном диоде Шоттки:
	\begin{equation}
		\frac{S_I}{I^2} \propto \Keff^{-2/3} \propto T^{1/3}.
		\label{eq:noise-pred}
	\end{equation}
	Используя операциональное определение $\Keff = I/I_0$ из
	уравнения (\ref{eq:keff-operational}), масштабирование шума можно переписать как
	\begin{equation}
		\frac{S_I}{I^2} = C \, \left(\frac{I}{I_0}\right)^{-2/3} + \text{const},
		\label{eq:noise-operational}
	\end{equation}
	где $C$ — материал-зависимая константа порядка единицы, а аддитивная
	константа учитывает неприводимый тепловой шум. Эта форма непосредственно
	доступна в стандартном эксперименте по спектроскопии шума: $I$ и $I_0$
	измеряются по $I$–$V$ характеристике, $S_I$ регистрируется как функция
	смещения при фиксированной температуре. Предсказание состоит в том, что
	график $S_I/I^2$ в зависимости от $(I/I_0)^{-2/3}$ должен быть линейным,
	причём наклон не зависит от температуры. Это жёсткий проверяемый тест,
	не требующий подгонки показателя.
	
	\section{Обсуждение: унификация показателей масштабирования}
	
	Появление показателя $3/2$ не уникально для диодов Шоттки, а указывает на
	более глубокий организационный принцип, о чём свидетельствует его
	независимое появление в универсальности роста поверхности Кардара–Паризи–Чжана
	(KPZ)~\cite{KPZ}. В KPZ динамический показатель $z = 3/2$ управляет
	фрактальным огрублением растущей границы раздела. YUCT даёт единую
	интерпретацию этих явлений. \textbf{Фронт роста может быть концептуализирован
		как $d_f = 3$-мерная координационная сеть, где кинетический показатель
		$\beta_{\mathrm{IV}} = d_f/2 = 3/2$ естественным образом возникает из
		нашего предположительного соотношения (\ref{eq:bIV-df}).}
	Таким образом, показатель KPZ и показатель Шоттки не просто аналогичны;
	они являются двумя кинетическими проявлениями одного и того же
	фундаментального закона ошибок YUCT $\varepsilon \propto \Keff^{-2/3}$ в
	совершенно разных физических реализациях. \textbf{Эта связь в настоящее
		время является гипотезой, требующей дальнейшей теоретической работы для
		непосредственной связи лагранжиана YUCT с уравнением KPZ.}
	
	Чтобы сделать предсказание KPZ более конкретным, мы можем явно
	отождествить $\Keff$ с долей насыщенных коррелированных мод в растущей
	границе. В режиме насыщения KPZ-границы латерального размера $L$ ширина
	$W(L,t)$ приближается к стационарному значению $W_{\text{sat}}(L) \simeq
	\sqrt{\langle h^2 \rangle} \propto L^{\alpha}$ с $\alpha = 1/2$.
	Амплитуда флуктуаций $\delta W = W(t) - W_{\text{sat}}$ контролируется
	количеством активных (ненасыщенных) мод. Мы предлагаем отождествить
	$\Keff$ с $1 - \nu_{\text{unsat}}/\nu_{\text{tot}}$, долей мод, которые
	перешли в коррелированную стационарную фазу. Тогда YUCT предсказывает
	\begin{equation}
		\frac{\delta W}{W_{\text{sat}}} \propto \Keff^{-2/3},
		\label{eq:kpz-prediction}
	\end{equation}
	где $\Keff$ можно оценить по временному затуханию ширины как
	$\Keff \approx 1 - (dW/dt)_{\text{norm}}$. Это предсказание может быть
	проверено в численных симуляциях и в экспериментах по кинетическому
	огрублению с высоким разрешением (например, электроосаждение,
	молекулярно-лучевая эпитаксия).
	
	Эта унификация предполагает иерархию универсальностей, проиллюстрированную
	на рис. \ref{fig:hierarchy}:
	\begin{itemize}
		\item \textbf{Первый порядок (YUCT):} Фундаментальный закон ошибок
		координации $\varepsilon = \kappa_c \alpha K_{\mathrm{eff}}^{-2/3}$
		действует на самом глубоком уровне, независимо от конкретной подложки.
		\item \textbf{Второй порядок (транспорт):} Для электронного переноса в
		2D-гетероструктурах Шоттки кинетический показатель $\beta_{\mathrm{IV}}$
		возникает из фрактальной размерности $d_f$ как $\beta_{\mathrm{IV}} = d_f/2$
		(гипотеза).
		\item \textbf{Третий порядок (рост):} В явлениях роста поверхности,
		принадлежащих классу KPZ, тот же кинетический показатель появляется как
		динамический критический показатель $z = 3/2$, отражая природу $d_f = 3$
		растущей границы (интерпретация).
	\end{itemize}
	Таким образом, показатель $3/2$ — не случайность, зависящая от материала,
	а геометрическая отметка $d_f = 3$-координационного процесса.
	
	\begin{figure}[htbp]
		\centering
		\begin{tikzpicture}[
			node distance=1.8cm,
			box/.style={rectangle, draw=blue!60, fill=blue!10, rounded corners, 
				text width=3.5cm, align=center, font=\small},
			arrow/.style={->, >=stealth, thick}
			]
			\node[box, fill=red!15, draw=red!50!black] (yuct) 
			{\textbf{Первый порядок (YUCT)}\\[2pt]
				$\varepsilon = \kappa_c\alpha K_{\mathrm{eff}}^{-2/3}$\\
				(фундаментальный закон ошибок)};
			\node[box, right=of yuct] (transport) 
			{\textbf{Второй порядок (транспорт)}\\[2pt]
				$\beta_{\mathrm{IV}} = d_f/2$\\
				(гетероструктуры Шоттки)};
			\node[box, right=of transport] (growth) 
			{\textbf{Третий порядок (рост)}\\[2pt]
				$z = 3/2$\\
				(KPZ-огрубление поверхности)};
			\draw[arrow] (yuct) -- (transport);
			\draw[arrow] (transport) -- (growth);
		\end{tikzpicture}
		\caption{Иерархия универсальностей. Закон ошибок координации YUCT
			(первый порядок) определяет эффективную фрактальную размерность $d_f$
			и кинетические показатели, наблюдаемые в электронном транспорте
			(второй порядок) и в росте поверхности (третий порядок). Пунктирные
			стрелки обозначают гипотезы.}
		\label{fig:hierarchy}
	\end{figure}
	
	Эта перспектива подтверждается растущими экспериментальными данными о том,
	что фрактальная размерность границы Шоттки непосредственно влияет на
	характеристики прибора~\cite{FractalSchottky}. В этих исследованиях
	фрактальная размерность является внешне заданным геометрическим свойством
	(шероховатость). В YUCT размерность $d_f = 3$ для идеальных поперечных
	контактов является \emph{эмерджентной} — она возникает из динамики самой
	координационной сети, даже когда геометрическая шероховатость стремится
	к нулю. Таким образом, внешне заданная $d_f$ модулирует уже существующую
	внутреннюю координационную сеть, что открывает многообещающий путь к
	созданию приборов на основе координации.
	
	\section{Универсальность в диапазоне 50 порядков величины}
	В таблице \ref{tab:universal-beta} собраны системы из различных приложений
	YUCT, для которых показатель масштабирования ошибок был эмпирически
	определён или предсказан как $\beta = 2/3$. Поперечная гетероструктура
	Шоттки и растущая граница KPZ добавлены как системы, для которых тот же
	показатель предсказывается на основе извлечённой фрактальной размерности
	$d_f = 3$. Таблица охватывает более 50 порядков величины по соответствующей
	энергетической (или пространственной) шкале.
	
	\begin{table}[H]
		\centering
		\caption{Системы, демонстрирующие или, как ожидается, демонстрирующие
			универсальный показатель ошибок YUCT $\beta = 2/3$ ($0.67$).}
		\label{tab:universal-beta}
		\begin{tabular}{@{}llll@{}}
			\toprule
			\textbf{Система} & \textbf{Наблюдаемая} & \textbf{Статус} & \textbf{Ссылка}\\
			\midrule
			Химические связи (873)      & $E \propto r^{-0.67}$               & Подтверждено  & Приложение C \\
			Ошибки репликации ДНК    & $\varepsilon \propto \Keff^{-0.67}$ & Подтверждено  & Приложение E \\
			Метаболическое масштабирование (простое)& $B \propto M^{0.67}$               & Подтверждено  & Приложение D, E.7 \\
			Волатильность ВВП от солнечной активности & $\sigma_{\mathrm{GDP}} \propto W^{-0.67}$ & Подтверждено & Приложение F \\
			Флуктуации температуры реликтового излучения    & $\delta T/T \propto \Keff^{-0.67}$ & Подтверждено  & Приложение M \\
			Гравитационные волны       & $\delta\tau/\tau \propto M^{-0.67}$& Подтверждено  & Приложение G, V \\
			\textbf{Поперечный Шоттки} & $S_I/I^2 \propto T^{1/3}$          & \textbf{Предсказано} & Настоящая работа \\
			\textbf{KPZ-граница}\textsuperscript{*} & Флуктуации ширины $\propto \Keff^{-2/3}$& \textbf{Предсказано} & Настоящая работа, \cite{KPZ}\\
			\bottomrule
		\end{tabular}
		
		\vspace{0.3cm}
		{\footnotesize \textsuperscript{*}Определение $K_{\mathrm{eff}}$ для растущей KPZ-границы должно быть сконструировано по аналогии со случаем транспорта, например, как доля коррелированных степеней свободы внутри корреляционного объёма. Это предсказание предполагает такую конструкцию и остаётся экспериментально непроверенным.}
	\end{table}
	
	\section{Заключение}
	Показатели масштабирования тока от температуры $\beta_{\mathrm{IV}} = 3/2$
	и $1$, о которых сообщили Ang и др., являются кинетическими проявлениями
	эффективной фрактальной размерности координационной сети, управляющей
	термоэлектронным переносом в 2D-гетероструктурах. Через соотношения YUCT
	$\beta_{\mathrm{IV}} = d_f/2$ (гипотеза) и $\beta_{\mathrm{err}} = 2/d_f$
	экспериментальные данные дают $d_f = 3$ для поперечных контактов, что
	предсказывает показатель ошибки координации $\beta_{\mathrm{err}} = 2/3$
	($\approx 0.67$). Это предсказание, проверяемое с помощью шумовой
	спектроскопии, расширяет область применимости универсального закона
	масштабирования YUCT в область твердотельной наноэлектроники. Те же
	фрактальная размерность и показатель естественным образом возникают в
	KPZ-росте поверхности, что указывает на глубокое геометрическое единство
	неравновесных явлений. YUCT даёт обоснование этому единству, хотя связи
	между теорией и как кинетическим соотношением Шоттки, так и уравнением
	KPZ остаются гипотезами, требующими дальнейшего исследования.
	Операциональные определения $\Keff$ для обоих контекстов были предложены
	для облегчения экспериментальных проверок.
	
	\begin{thebibliography}{99}
		\bibitem{Ang2018} Y.~S.~Ang, H.~Y.~Yang, and L.~K.~Ang,
		Phys. Rev. Lett. \textbf{121}, 056802 (2018).
		\bibitem{KPZ} M.~Kardar, G.~Parisi, Y.-C.~Zhang,
		Phys. Rev. Lett. \textbf{56}, 889 (1986).
		\bibitem{FractalSchottky} S.~Zeybek, M.~Biber, A.~Türüt,
		Appl. Surf. Sci. \textbf{253}, 3881 (2007). [Представительное исследование влияния фрактальной размерности на диоды Шоттки.]
		\bibitem{Yakushev2026} A.~V.~Yakushev, \emph{Yakushev Unified Coordination Theory (YUCT)}, Zenodo (2026).
		\bibitem{AppendixL} A.~V.~Yakushev, ``Приложение L: Фрактальное масштабирование координационной ошибки'', в \emph{YUCT} (2026).
		\bibitem{AppendixM} A.~V.~Yakushev, ``Приложение M: Космология YUCT'', в \emph{YUCT} (2026).
		\bibitem{AppendixE} A.~V.~Yakushev, ``Приложение E: Координационная генетика'', в \emph{YUCT} (2026).
		\bibitem{AppendixF} A.~V.~Yakushev, ``Приложение F: Применение YUCT к экономике'', в \emph{YUCT} (2026).
		\bibitem{AppendixG} A.~V.~Yakushev, ``Приложение G: Решение проблемы квантовой гравитации'', в \emph{YUCT} (2026).
		\bibitem{AppendixV} A.~V.~Yakushev, ``Приложение V: Время как координационная энтропия'', в \emph{YUCT} (2026).
	\end{thebibliography}
	
\end{document}